\documentclass[11pt,a4paper]{article}
\usepackage[utf8]{inputenc}
\usepackage[T1]{fontenc}
\usepackage[catalan]{babel}
\usepackage{lmodern}
\usepackage[sfdefault,lining]{FiraSans}
\usepackage[a4paper,top=2.3cm,bottom=2.2cm,left=1.8cm,right=1.8cm,headheight=1.3cm,headsep=0.4cm,footskip=1cm]{geometry}
\usepackage{amsmath,amssymb}
\usepackage{icomma}
\usepackage{graphicx}
\usepackage[table]{xcolor}
\usepackage{tikz}
\usepackage{fontawesome5}
\usepackage{ccicons}
\usepackage{fancyhdr}
\usepackage{lastpage}
\usepackage{tasks}
\usepackage{needspace}
\usepackage[most]{tcolorbox}
\usepackage{array}
\usepackage[hidelinks,pdftitle={Fitxa d'exercicis de la unitat 9: Estadística Unidimensional. 4t d'ESO Matemàtiques A},pdfsubject={Matemàtiques A 4t ESO},pdfkeywords={estadística, població, mostra, mostreig, variables, taules de freqüències, dades agrupades, gràfics estadístics, moda, mitjana, mediana, rang, desviació mitjana, variància, desviació típica, coeficient de variació, quartils, percentils, diagrama de caixes}]{hyperref}

\definecolor{navy}{HTML}{1B2748}
\definecolor{accent}{HTML}{E8603A}
\definecolor{teal}{HTML}{1F9E90}
\definecolor{soft}{HTML}{EEF1F8}
\definecolor{muted}{HTML}{5D6685}

\setlength{\parindent}{0pt}
\setlength{\parskip}{3pt}
\settasks{label=\alph*), label-format=\bfseries\color{navy}, label-width=1.3em, item-indent=1.8em, column-sep=0.6em, after-item-skip=5pt, before-skip=3pt, after-skip=4pt}

% ---------- capçalera i peu ----------
\newcommand{\llicencia}{\href{https://creativecommons.org/licenses/by-nc-sa/4.0/deed.ca}{\ccbyncsa}\ \ Llicència Creative Commons Reconeixement-NoComercial-CompartirIgual 4.0 Internacional (CC BY-NC-SA 4.0)}
\pagestyle{fancy}
\fancyhf{}
\fancyhead[L]{\raisebox{-0.35\height}{\includegraphics[height=0.95cm]{logo-gva.png}}}
\fancyhead[C]{\small\color{muted}Fitxa d'exercicis · Unitat 9: Estadística Unidimensional · 4t ESO A}
\fancyhead[R]{\raisebox{-0.35\height}{\includegraphics[height=0.95cm]{logo-vermellar.png}}}
\fancyfoot[L]{\scriptsize\color{muted}\llicencia}
\fancyfoot[R]{\small\color{muted}\thepage}
\renewcommand{\headrulewidth}{0.4pt}
\renewcommand{\headrule}{\color{muted!40}\hrule height \headrulewidth}
\fancypagestyle{primera}{\fancyhf{}\renewcommand{\headrulewidth}{0pt}
  \fancyfoot[L]{\scriptsize\color{muted}\llicencia}
  \fancyfoot[R]{\small\color{muted}\thepage}}

% ---------- nivells i exercicis ----------
\newcommand{\estrelles}[1]{{\color{accent}\ifcase#1\or\faStar\faStar[regular]\faStar[regular]\or\faStar\faStar\faStar[regular]\or\faStar\faStar\faStar\fi}}
\newcounter{exer}
\newcommand{\ex}[1]{\par\needspace{4\baselineskip}\medskip\stepcounter{exer}\noindent{\color{navy}\bfseries\theexer.}\ \ {\small\estrelles{#1}}\ \ }

\newcommand{\seccio}[2]{\par\needspace{10\baselineskip}\bigskip
  \begin{tcolorbox}[colback=navy,colframe=navy,coltext=white,arc=2mm,boxrule=0pt,left=3mm,right=3mm,top=1.2mm,bottom=1.2mm]
  {\large\bfseries #1\quad #2}
  \end{tcolorbox}\vspace{-2pt}}

\newcolumntype{C}[1]{>{\centering\arraybackslash}p{#1}}
\definecolor{violet}{HTML}{7C6CF0}
\definecolor{gold}{HTML}{E0A21B}
% quadre de recordatori (només a la fitxa d'exercicis)
\newcommand{\recorda}[1]{\par\smallskip\begin{tcolorbox}[colback=accent!7,colframe=accent!55,arc=2mm,boxrule=0.5pt,left=2.5mm,right=2.5mm,top=1mm,bottom=1mm]\small{\color{accent}\faLightbulb[regular]}\ \textbf{Recorda:}\ #1\end{tcolorbox}}
\newcommand{\res}[1]{#1}\newcommand{\rest}[1]{#1}

\def\fmneed{14}
% =====================================================================
% Capçalera FentMates de les fitxes: un color per a cada curs.
% Els colors són els de la pàgina inicial (LEVELS a assets/js/materials.js);
% el logo ve de _eines/logo/pdf/ (python3 _eines/logo/genera.py).
%
% Ús, al final del preàmbul de la unitat:
%   % =====================================================================
% Capçalera FentMates de les fitxes: un color per a cada curs.
% Els colors són els de la pàgina inicial (LEVELS a assets/js/materials.js);
% el logo ve de _eines/logo/pdf/ (python3 _eines/logo/genera.py).
%
% Ús, al final del preàmbul de la unitat:
%   % =====================================================================
% Capçalera FentMates de les fitxes: un color per a cada curs.
% Els colors són els de la pàgina inicial (LEVELS a assets/js/materials.js);
% el logo ve de _eines/logo/pdf/ (python3 _eines/logo/genera.py).
%
% Ús, al final del preàmbul de la unitat:
%   \input{../comu/fentmates.tex}
%   \fmcurs{1eso}
% i, al començament del document:
%   \capcalera{Matemàtiques}{1r ESO}{Fitxa d'exercicis · Unitat 2}{Divisibilitat}
% Opcional abans de \input: \def\fmneed{14} (espai mínim abans de cada secció).
% =====================================================================
\makeatletter
\def\fm@logodir{../../logo/pdf}
\providecommand{\fmneed}{10}
\def\fm@cursnom{}

% clau del curs, color i color del text damunt del color
\newcommand{\fm@curs}[3]{\@namedef{fm@color@#1}{#2}\@namedef{fm@text@#1}{#3}}
\fm@curs{1eso}{E6B84A}{navy}
\fm@curs{2eso}{3B82F6}{white}
\fm@curs{3eso}{FF7A50}{white}
\fm@curs{4eso-a}{A78BFA}{white}
\fm@curs{4eso-b}{6D28D9}{white}
\fm@curs{1batx-mi}{9B1C1C}{white}
\fm@curs{1batx-ccss}{EF6B6B}{white}
\fm@curs{2batx-mii}{1E6B47}{white}
\fm@curs{2batx-ccss}{2BB3A3}{white}

\newcommand{\fmcurs}[1]{%
  \@ifundefined{fm@color@#1}{\PackageError{fentmates}{Curs desconegut: #1}{}}{}%
  \edef\fm@tmp{\noexpand\definecolor{curs}{HTML}{\@nameuse{fm@color@#1}}%
               \noexpand\colorlet{curstext}{\@nameuse{fm@text@#1}}}\fm@tmp
  \colorlet{cursfosc}{curs!72!black}%
  \def\fm@logo{\fm@logodir/vertical-#1.pdf}%
  \def\fm@marca{\fm@logodir/marca-#1.pdf}}

% ---------- capçalera de la primera pàgina ----------
% #1 matèria (petit), #2 curs, #3 tipus de document i unitat, #4 títol
\newcommand{\capcalera}[4]{%
  \gdef\fm@cursnom{#2}%
  \par\noindent
  \begin{tcolorbox}[colback=curs,colframe=curs,boxrule=0pt,arc=2.2mm,width=3.5cm,height=2.15cm,
      valign=center,left=3mm,right=2mm,top=1mm,bottom=1mm,nobeforeafter,box align=center]
    \color{curstext}{\footnotesize #1}\par\vspace{1pt}{\fontsize{19}{21}\selectfont\bfseries #2\par}%
  \end{tcolorbox}\hspace{5mm}%
  \begin{minipage}[c]{\dimexpr\linewidth-3.5cm-5mm-2.9cm\relax}
    {\color{cursfosc}\small\bfseries\MakeUppercase{#3}}\par\vspace{2pt}
    {\color{navy}\fontsize{20}{23}\selectfont\bfseries #4\par}
  \end{minipage}\hfill
  \begin{minipage}[c]{2.6cm}\raggedleft\includegraphics[width=2.6cm]{\fm@logo}\end{minipage}\par
  \vspace{2pt}}

% ---------- barres de secció del color del curs ----------
\renewcommand{\seccio}[2]{\par\needspace{\fmneed\baselineskip}\bigskip
  \begin{tcolorbox}[colback=curs,colframe=curs,coltext=curstext,arc=2mm,boxrule=0pt,left=3mm,right=3mm,top=1.2mm,bottom=1.2mm]
  {\large\bfseries #1\quad #2}
  \end{tcolorbox}\vspace{-2pt}}

% ---------- estrelles de dificultat del color del curs (a Batxillerat no se'n posen) ----------
\@ifundefined{estrelles}{}{%
  \renewcommand{\estrelles}[1]{{\color{curs}\ifcase#1\or\faStar\faStar[regular]\faStar[regular]\or\faStar\faStar\faStar[regular]\or\faStar\faStar\faStar\fi}}}

% ---------- capçalera i peu de les pàgines ----------
% primera pàgina: logos institucionals al peu
\fancypagestyle{primera}{\fancyhf{}\renewcommand{\headrulewidth}{0pt}
  \fancyfoot[L]{\raisebox{-0.32\height}{\includegraphics[height=0.78cm]{logo-gva.png}}\hspace{5mm}%
    \raisebox{-0.32\height}{\includegraphics[height=0.78cm]{logo-vermellar.png}}}
  \fancyfoot[R]{\scriptsize\color{muted}\href{https://creativecommons.org/licenses/by-nc-sa/4.0/deed.ca}{\ccbyncsa}\ \ CC BY-NC-SA 4.0 · fentmates.org\qquad\small\thepage}}
% la resta: el símbol a l'esquerra i el curs a la dreta, amb el títol al centre com abans
\fancyhead[L]{\raisebox{-0.3\height}{\includegraphics[height=0.62cm]{\fm@marca}}}
\fancyhead[R]{\small\bfseries\color{cursfosc}\fm@cursnom}
\renewcommand{\headrule}{\color{curs!45}\hrule height \headrulewidth}
\makeatother

%   \fmcurs{1eso}
% i, al començament del document:
%   \capcalera{Matemàtiques}{1r ESO}{Fitxa d'exercicis · Unitat 2}{Divisibilitat}
% Opcional abans de \input: \def\fmneed{14} (espai mínim abans de cada secció).
% =====================================================================
\makeatletter
\def\fm@logodir{../../logo/pdf}
\providecommand{\fmneed}{10}
\def\fm@cursnom{}

% clau del curs, color i color del text damunt del color
\newcommand{\fm@curs}[3]{\@namedef{fm@color@#1}{#2}\@namedef{fm@text@#1}{#3}}
\fm@curs{1eso}{E6B84A}{navy}
\fm@curs{2eso}{3B82F6}{white}
\fm@curs{3eso}{FF7A50}{white}
\fm@curs{4eso-a}{A78BFA}{white}
\fm@curs{4eso-b}{6D28D9}{white}
\fm@curs{1batx-mi}{9B1C1C}{white}
\fm@curs{1batx-ccss}{EF6B6B}{white}
\fm@curs{2batx-mii}{1E6B47}{white}
\fm@curs{2batx-ccss}{2BB3A3}{white}

\newcommand{\fmcurs}[1]{%
  \@ifundefined{fm@color@#1}{\PackageError{fentmates}{Curs desconegut: #1}{}}{}%
  \edef\fm@tmp{\noexpand\definecolor{curs}{HTML}{\@nameuse{fm@color@#1}}%
               \noexpand\colorlet{curstext}{\@nameuse{fm@text@#1}}}\fm@tmp
  \colorlet{cursfosc}{curs!72!black}%
  \def\fm@logo{\fm@logodir/vertical-#1.pdf}%
  \def\fm@marca{\fm@logodir/marca-#1.pdf}}

% ---------- capçalera de la primera pàgina ----------
% #1 matèria (petit), #2 curs, #3 tipus de document i unitat, #4 títol
\newcommand{\capcalera}[4]{%
  \gdef\fm@cursnom{#2}%
  \par\noindent
  \begin{tcolorbox}[colback=curs,colframe=curs,boxrule=0pt,arc=2.2mm,width=3.5cm,height=2.15cm,
      valign=center,left=3mm,right=2mm,top=1mm,bottom=1mm,nobeforeafter,box align=center]
    \color{curstext}{\footnotesize #1}\par\vspace{1pt}{\fontsize{19}{21}\selectfont\bfseries #2\par}%
  \end{tcolorbox}\hspace{5mm}%
  \begin{minipage}[c]{\dimexpr\linewidth-3.5cm-5mm-2.9cm\relax}
    {\color{cursfosc}\small\bfseries\MakeUppercase{#3}}\par\vspace{2pt}
    {\color{navy}\fontsize{20}{23}\selectfont\bfseries #4\par}
  \end{minipage}\hfill
  \begin{minipage}[c]{2.6cm}\raggedleft\includegraphics[width=2.6cm]{\fm@logo}\end{minipage}\par
  \vspace{2pt}}

% ---------- barres de secció del color del curs ----------
\renewcommand{\seccio}[2]{\par\needspace{\fmneed\baselineskip}\bigskip
  \begin{tcolorbox}[colback=curs,colframe=curs,coltext=curstext,arc=2mm,boxrule=0pt,left=3mm,right=3mm,top=1.2mm,bottom=1.2mm]
  {\large\bfseries #1\quad #2}
  \end{tcolorbox}\vspace{-2pt}}

% ---------- estrelles de dificultat del color del curs (a Batxillerat no se'n posen) ----------
\@ifundefined{estrelles}{}{%
  \renewcommand{\estrelles}[1]{{\color{curs}\ifcase#1\or\faStar\faStar[regular]\faStar[regular]\or\faStar\faStar\faStar[regular]\or\faStar\faStar\faStar\fi}}}

% ---------- capçalera i peu de les pàgines ----------
% primera pàgina: logos institucionals al peu
\fancypagestyle{primera}{\fancyhf{}\renewcommand{\headrulewidth}{0pt}
  \fancyfoot[L]{\raisebox{-0.32\height}{\includegraphics[height=0.78cm]{logo-gva.png}}\hspace{5mm}%
    \raisebox{-0.32\height}{\includegraphics[height=0.78cm]{logo-vermellar.png}}}
  \fancyfoot[R]{\scriptsize\color{muted}\href{https://creativecommons.org/licenses/by-nc-sa/4.0/deed.ca}{\ccbyncsa}\ \ CC BY-NC-SA 4.0 · fentmates.org\qquad\small\thepage}}
% la resta: el símbol a l'esquerra i el curs a la dreta, amb el títol al centre com abans
\fancyhead[L]{\raisebox{-0.3\height}{\includegraphics[height=0.62cm]{\fm@marca}}}
\fancyhead[R]{\small\bfseries\color{cursfosc}\fm@cursnom}
\renewcommand{\headrule}{\color{curs!45}\hrule height \headrulewidth}
\makeatother

%   \fmcurs{1eso}
% i, al començament del document:
%   \capcalera{Matemàtiques}{1r ESO}{Fitxa d'exercicis · Unitat 2}{Divisibilitat}
% Opcional abans de \input: \def\fmneed{14} (espai mínim abans de cada secció).
% =====================================================================
\makeatletter
\def\fm@logodir{../../logo/pdf}
\providecommand{\fmneed}{10}
\def\fm@cursnom{}

% clau del curs, color i color del text damunt del color
\newcommand{\fm@curs}[3]{\@namedef{fm@color@#1}{#2}\@namedef{fm@text@#1}{#3}}
\fm@curs{1eso}{E6B84A}{navy}
\fm@curs{2eso}{3B82F6}{white}
\fm@curs{3eso}{FF7A50}{white}
\fm@curs{4eso-a}{A78BFA}{white}
\fm@curs{4eso-b}{6D28D9}{white}
\fm@curs{1batx-mi}{9B1C1C}{white}
\fm@curs{1batx-ccss}{EF6B6B}{white}
\fm@curs{2batx-mii}{1E6B47}{white}
\fm@curs{2batx-ccss}{2BB3A3}{white}

\newcommand{\fmcurs}[1]{%
  \@ifundefined{fm@color@#1}{\PackageError{fentmates}{Curs desconegut: #1}{}}{}%
  \edef\fm@tmp{\noexpand\definecolor{curs}{HTML}{\@nameuse{fm@color@#1}}%
               \noexpand\colorlet{curstext}{\@nameuse{fm@text@#1}}}\fm@tmp
  \colorlet{cursfosc}{curs!72!black}%
  \def\fm@logo{\fm@logodir/vertical-#1.pdf}%
  \def\fm@marca{\fm@logodir/marca-#1.pdf}}

% ---------- capçalera de la primera pàgina ----------
% #1 matèria (petit), #2 curs, #3 tipus de document i unitat, #4 títol
\newcommand{\capcalera}[4]{%
  \gdef\fm@cursnom{#2}%
  \par\noindent
  \begin{tcolorbox}[colback=curs,colframe=curs,boxrule=0pt,arc=2.2mm,width=3.5cm,height=2.15cm,
      valign=center,left=3mm,right=2mm,top=1mm,bottom=1mm,nobeforeafter,box align=center]
    \color{curstext}{\footnotesize #1}\par\vspace{1pt}{\fontsize{19}{21}\selectfont\bfseries #2\par}%
  \end{tcolorbox}\hspace{5mm}%
  \begin{minipage}[c]{\dimexpr\linewidth-3.5cm-5mm-2.9cm\relax}
    {\color{cursfosc}\small\bfseries\MakeUppercase{#3}}\par\vspace{2pt}
    {\color{navy}\fontsize{20}{23}\selectfont\bfseries #4\par}
  \end{minipage}\hfill
  \begin{minipage}[c]{2.6cm}\raggedleft\includegraphics[width=2.6cm]{\fm@logo}\end{minipage}\par
  \vspace{2pt}}

% ---------- barres de secció del color del curs ----------
\renewcommand{\seccio}[2]{\par\needspace{\fmneed\baselineskip}\bigskip
  \begin{tcolorbox}[colback=curs,colframe=curs,coltext=curstext,arc=2mm,boxrule=0pt,left=3mm,right=3mm,top=1.2mm,bottom=1.2mm]
  {\large\bfseries #1\quad #2}
  \end{tcolorbox}\vspace{-2pt}}

% ---------- estrelles de dificultat del color del curs (a Batxillerat no se'n posen) ----------
\@ifundefined{estrelles}{}{%
  \renewcommand{\estrelles}[1]{{\color{curs}\ifcase#1\or\faStar\faStar[regular]\faStar[regular]\or\faStar\faStar\faStar[regular]\or\faStar\faStar\faStar\fi}}}

% ---------- capçalera i peu de les pàgines ----------
% primera pàgina: logos institucionals al peu
\fancypagestyle{primera}{\fancyhf{}\renewcommand{\headrulewidth}{0pt}
  \fancyfoot[L]{\raisebox{-0.32\height}{\includegraphics[height=0.78cm]{logo-gva.png}}\hspace{5mm}%
    \raisebox{-0.32\height}{\includegraphics[height=0.78cm]{logo-vermellar.png}}}
  \fancyfoot[R]{\scriptsize\color{muted}\href{https://creativecommons.org/licenses/by-nc-sa/4.0/deed.ca}{\ccbyncsa}\ \ CC BY-NC-SA 4.0 · fentmates.org\qquad\small\thepage}}
% la resta: el símbol a l'esquerra i el curs a la dreta, amb el títol al centre com abans
\fancyhead[L]{\raisebox{-0.3\height}{\includegraphics[height=0.62cm]{\fm@marca}}}
\fancyhead[R]{\small\bfseries\color{cursfosc}\fm@cursnom}
\renewcommand{\headrule}{\color{curs!45}\hrule height \headrulewidth}
\makeatother

\fmcurs{4eso-a}

\begin{document}
\thispagestyle{primera}

\capcalera{Matemàtiques A}{4t ESO}{Fitxa d'exercicis · Unitat 9}{Estadística Unidimensional}
\vspace{10pt}
Nom i cognoms: \rule{8.4cm}{0.4pt}\hfill Grup: \rule{1.6cm}{0.4pt}\hfill Data: \rule{2.2cm}{0.4pt}
\vspace{8pt}

\begin{tcolorbox}[colback=soft,colframe=soft,arc=2mm,left=3mm,right=3mm,top=1.5mm,bottom=1.5mm]
\small Exercicis graduats per nivells de dificultat:\quad \estrelles{1}\ bàsic \qquad \estrelles{2}\ mitjà \qquad \estrelles{3}\ avançat.\\[2pt] Resol-los al quadern mostrant tots els passos. Pots fer servir la calculadora en mode estadístic per a comprovar els resultats. Les seccions segueixen l'ordre dels apunts de la unitat (9.1--9.8).
\end{tcolorbox}

% =====================================================
\seccio{1}{Població, mostra i variables estadístiques}

\ex{1} Classifica aquestes variables en qualitatives, quantitatives discretes o quantitatives contínues:\begin{tasks}(2)
  \task nombre de missatges enviats en un dia
  \task pes d'una maleta
  \task marca del cotxe familiar
  \task temps d'espera a la parada de l'autobús
  \task nombre de germans
  \task nivell d'estudis (ESO, Batxillerat…)
\end{tasks}\recorda{Si els valors no són nombres, és \textbf{qualitativa}. Si són nombres que es \textbf{compten}, és discreta; si es \textbf{mesuren}, és contínua.}

\ex{1} Per a saber quantes hores d'esport fan a la setmana els 900 alumnes d'un institut, se n'enquesta 90 triats a l'atzar. Indica la població, la mostra i la seua grandària, un individu i la variable estudiada (i de quin tipus és).

\ex{2} Raona si aquestes mostres són representatives:\begin{tasks}(1)
  \task Per a conéixer l'esport preferit dels joves d'una ciutat, es pregunta a 50 persones que ixen d'un camp de futbol.
  \task Per a conéixer l'alçada mitjana de l'alumnat d'un institut, es trien a l'atzar 5 alumnes de cada grup.
\end{tasks}

\ex{2} Un poble té 2\,400 habitants: 600 menors de 18 anys, 1\,200 d'entre 18 i 65 anys i 600 de més de 65. Es vol fer una enquesta a 120 persones amb un mostreig estratificat. Quantes persones s'han de triar de cada grup d'edat?\recorda{En el mostreig \textbf{estratificat}, cada grup aporta a la mostra la mateixa proporció que té en la població.}

\ex{3} De la llista dels 840 alumnes d'un institut volem triar-ne 60 amb un mostreig sistemàtic. Cada quants alumnes n'hem d'agafar un? Si el primer triat a l'atzar és el número 9, quins són el segon, el cinqué i l'últim?


% =====================================================
\seccio{2}{Taules de freqüències amb dades aïllades}

\ex{1} Aquests són els llibres que han llegit en un mes 24 alumnes: \ $1,\ \allowbreak 2,\ \allowbreak 4,\ \allowbreak 2,\ \allowbreak 0,\ \allowbreak 1,\ \allowbreak 3,\ \allowbreak 2,\ \allowbreak 2,\ \allowbreak 2,\ \allowbreak 5,\ \allowbreak 2,\ \allowbreak 1,\ \allowbreak 1,\ \allowbreak 3,\ \allowbreak 4,\ \allowbreak 0,\ \allowbreak 0,\ \allowbreak 2,\ \allowbreak 2,\ \allowbreak 1,\ \allowbreak 3,\ \allowbreak 3,\ \allowbreak 1$. \ Fes la taula de freqüències amb $f_i$, $F_i$, $h_i$ i percentatges.

\ex{1} Amb la taula de l'exercici anterior, respon:\begin{tasks}(2)
  \task Quants alumnes han llegit com a molt 2 llibres?
  \task Quants n'han llegit més de 3?
  \task Quin percentatge n'ha llegit 3 o menys?
  \task Quin percentatge no n'ha llegit cap?
\end{tasks}

\ex{2} Completa la taula de l'esport preferit de 40 alumnes:\par\smallskip\begin{center}{\renewcommand{\arraystretch}{1.25}\begin{tabular}{|c|C{1.2cm}|C{1.2cm}|C{1.2cm}|}\hline Esport & $f_i$ & $h_i$ & \%\\\hline Futbol & 14 & & \\\hline Bàsquet & & $0{,}25$ & \\\hline Tenis & & & 15\\\hline Natació & 8 & & \\\hline Altres & & & \\\hline \textbf{Total} & 40 & & \\\hline\end{tabular}}\end{center}

\ex{2} En aquesta taula, $x_i$ és el nombre de visites al dentista en un any i $F_i$, la freqüència acumulada. Troba les freqüències absolutes. Quantes persones hi han anat exactament 3 vegades? I almenys 2 vegades?\par\smallskip\begin{center}{\renewcommand{\arraystretch}{1.3}\begin{tabular}{|c|C{0.9cm}|C{0.9cm}|C{0.9cm}|C{0.9cm}|C{0.9cm}|C{0.9cm}|}\hline $x_i$ & $0$ & $1$ & $2$ & $3$ & $4$ & $5$\\\hline $F_i$ & $2$ & $7$ & $15$ & $22$ & $26$ & $28$\\\hline\end{tabular}}\end{center}

\ex{2} Hem llançat un dau 60 vegades. Ha eixit l'1 nou vegades, el 2 onze, el 3 deu, el 4 huit, el 5 dotze i el 6 la resta. Quantes vegades ha eixit el 6? Quin percentatge dels llançaments ha eixit un nombre parell? I un nombre més gran que 4?

\ex{3} Una variable pren tres valors, $A$, $B$ i $C$. La freqüència relativa d'$A$ és $0{,}25$; la de $B$, $0{,}4$, i la freqüència absoluta de $B$ és 16. Calcula $N$, la freqüència relativa de $C$ i totes les freqüències absolutes.

\ex{3} En una taula de freqüències, $F_3=12$, $F_4=19$ i $F_5=25$, i 5 és el valor més gran. Quantes dades hi ha? Quina és la freqüència absoluta del 4? Quin percentatge de les dades és més gran que 3?


% =====================================================
\seccio{3}{Taules amb dades agrupades en intervals}

\ex{1} Els intervals d'una taula són $[10,\ 15)$, $[15,\ 20)$, $[20,\ 25)$ i $[25,\ 30)$. En quin interval va cadascuna d'aquestes dades? Quina és la marca de classe de cada interval?\begin{tasks}(4)
  \task $15$
  \task $24{,}9$
  \task $25$
  \task $10$
\end{tasks}

\ex{1} Aquests són els pesos (en kg) de 24 alumnes: \ $48,\ \allowbreak 55,\ \allowbreak 61,\ \allowbreak 44,\ \allowbreak 52,\ \allowbreak 58,\ \allowbreak 63,\ \allowbreak 47,\ \allowbreak 50,\ \allowbreak 56,\ \allowbreak 59,\ \allowbreak 66,\ \allowbreak 45,\ \allowbreak 53,\ \allowbreak 57,\ \allowbreak 60,\ \allowbreak 49,\ \allowbreak 54,\ \allowbreak 62,\ \allowbreak 51,\ \allowbreak 58,\ \allowbreak 64,\ \allowbreak 46,\ \allowbreak 55$. \ Agrupa'ls en intervals d'amplitud 5 des de 40\,kg i fes la taula amb les marques de classe, $f_i$, $F_i$ i percentatges. Quin percentatge pesa menys de 55\,kg?\recorda{L'interval $[a,\ b)$ conté les dades $\geq a$ i $<b$. La \textbf{marca de classe} és el punt mitjà: $\frac{a+b}{2}$.}

\ex{2} Les puntuacions de 40 alumnes en un test han sigut: \ $52,\ \allowbreak 67,\ \allowbreak 45,\ \allowbreak 73,\ \allowbreak 38,\ \allowbreak 61,\ \allowbreak 55,\ \allowbreak 80,\ \allowbreak 49,\ \allowbreak 64,\ \allowbreak 58,\ \allowbreak 71,\ \allowbreak 43,\ \allowbreak 66,\ \allowbreak 59,\ \allowbreak 35,\ \allowbreak 77,\ \allowbreak 62,\ \allowbreak 50,\ \allowbreak 69,\ \allowbreak 57,\ \allowbreak 84,\ \allowbreak 47,\ \allowbreak 63,\ \allowbreak 54,\ \allowbreak 74,\ \allowbreak 41,\ \allowbreak 68,\ \allowbreak 60,\ \allowbreak 56,\ \allowbreak 31,\ \allowbreak 79,\ \allowbreak 65,\ \allowbreak 51,\ \allowbreak 88,\ \allowbreak 70,\ \allowbreak 46,\ \allowbreak 58,\ \allowbreak 62,\ \allowbreak 53$. \ Decideix quants intervals fas i de quina amplitud (fes servir la regla de l'arrel) i fes la taula de freqüències. Quin percentatge ha tret menys de 50 punts?\recorda{Rang $=x_{\max}-x_{\min}$. \ Nombre d'intervals $\approx\sqrt N$. \ Amplitud $\approx\frac{\text{rang}}{\text{nre. d'intervals}}$, arrodonida cap amunt.}

\ex{2} La taula recull les alçades (en cm) de 50 alumnes. Fes la columna de les $F_i$ i respon: quants alumnes fan menys de 165\,cm? Quin percentatge fa 170\,cm o més? Quants fan entre 155 i 170\,cm?\par\smallskip\begin{center}{\renewcommand{\arraystretch}{1.3}\begin{tabular}{|c|C{1.6cm}|C{1.6cm}|C{1.6cm}|C{1.6cm}|C{1.6cm}|C{1.6cm}|}\hline Alçada & $[150,\ 155)$ & $[155,\ 160)$ & $[160,\ 165)$ & $[165,\ 170)$ & $[170,\ 175)$ & $[175,\ 180)$\\\hline $f_i$ & $4$ & $9$ & $15$ & $12$ & $7$ & $3$\\\hline\end{tabular}}\end{center}

\ex{3} Les marques de classe d'una taula són 4, 8, 12 i 16, i tots els intervals tenen la mateixa amplitud. Quins són els intervals? I si un interval té marca $32{,}5$ i amplitud 5?


% =====================================================
\seccio{4}{Gràfics estadístics}

\ex{1} Dibuixa el diagrama de barres i el polígon de freqüències dels llibres llegits pels 24 alumnes de la secció 2 (exercici 6).\recorda{En un \textbf{diagrama de barres}, les barres van separades i tenen l'altura de la freqüència. El \textbf{polígon} uneix els extrems de dalt de les barres.}

\ex{1} Els 36 alumnes d'una classe vénen a l'institut així: a peu 15, en autobús 12, en cotxe 6 i en bici 3. Calcula l'angle de cada sector i dibuixa el diagrama de sectors.\recorda{Angle del sector $=\dfrac{f_i}{N}\cdot360^\circ$.}

\ex{2} Aquest diagrama mostra quantes vegades han viatjat en avió els alumnes d'una classe. Fes la taula de freqüències. Quants alumnes hi ha? Quina és la moda? Quants han viatjat 3 o més vegades?\par\smallskip\begin{center}\begin{tikzpicture}[baseline=(current bounding box.center),line cap=round,line join=round]
\draw[muted!30,line width=0.3pt] (0,0.000)--(4.700,0.000); \node[left,font=\tiny,text=muted,inner sep=1.5pt] at (0,0.000) {$0$};
\draw[muted!30,line width=0.3pt] (0,0.350)--(4.700,0.350); \node[left,font=\tiny,text=muted,inner sep=1.5pt] at (0,0.350) {$1$};
\draw[muted!30,line width=0.3pt] (0,0.700)--(4.700,0.700); \node[left,font=\tiny,text=muted,inner sep=1.5pt] at (0,0.700) {$2$};
\draw[muted!30,line width=0.3pt] (0,1.050)--(4.700,1.050); \node[left,font=\tiny,text=muted,inner sep=1.5pt] at (0,1.050) {$3$};
\draw[muted!30,line width=0.3pt] (0,1.400)--(4.700,1.400); \node[left,font=\tiny,text=muted,inner sep=1.5pt] at (0,1.400) {$4$};
\draw[muted!30,line width=0.3pt] (0,1.750)--(4.700,1.750); \node[left,font=\tiny,text=muted,inner sep=1.5pt] at (0,1.750) {$5$};
\draw[muted!30,line width=0.3pt] (0,2.100)--(4.700,2.100); \node[left,font=\tiny,text=muted,inner sep=1.5pt] at (0,2.100) {$6$};
\draw[muted!30,line width=0.3pt] (0,2.450)--(4.700,2.450); \node[left,font=\tiny,text=muted,inner sep=1.5pt] at (0,2.450) {$7$};
\draw[muted!30,line width=0.3pt] (0,2.800)--(4.700,2.800); \node[left,font=\tiny,text=muted,inner sep=1.5pt] at (0,2.800) {$8$};
\draw[-latex,muted,line width=0.6pt] (0,0)--(5.000,0) node[right,font=\scriptsize,text=navy]{Viatges};
\draw[-latex,muted,line width=0.6pt] (0,0)--(0,3.150) node[above,font=\scriptsize,text=navy]{Alumnes};
\draw[teal,fill=teal!35] (0.280,0) rectangle (0.820,1.750);
\node[below,font=\scriptsize,inner sep=2pt] at (0.550,0) {0};
\draw[teal,fill=teal!35] (1.180,0) rectangle (1.720,2.800);
\node[below,font=\scriptsize,inner sep=2pt] at (1.450,0) {1};
\draw[teal,fill=teal!35] (2.080,0) rectangle (2.620,2.100);
\node[below,font=\scriptsize,inner sep=2pt] at (2.350,0) {2};
\draw[teal,fill=teal!35] (2.980,0) rectangle (3.520,1.400);
\node[below,font=\scriptsize,inner sep=2pt] at (3.250,0) {3};
\draw[teal,fill=teal!35] (3.880,0) rectangle (4.420,0.700);
\node[below,font=\scriptsize,inner sep=2pt] at (4.150,0) {4};
\end{tikzpicture}\end{center}

\ex{2} Dibuixa l'histograma i el polígon de freqüències de les alçades dels 50 alumnes de la secció 3 (exercici 16).\recorda{En un \textbf{histograma} les barres ocupen tot l'interval i van juntes. El polígon uneix els punts de dalt sobre les marques de classe i es tanca en l'eix.}

\ex{2} S'ha preguntat a 120 persones on prefereixen passar les vacances. Quantes persones han triat cada lloc?\par\smallskip\begin{center}\begin{tikzpicture}[baseline=(current bounding box.center)]
\fill[teal!70,draw=white,line width=0.8pt] (0,0)--(90.000:1.45) arc[start angle=90.000,end angle=-18.000,radius=1.45cm]--cycle;
\node[font=\scriptsize\bfseries,text=white] at (36.000:0.899) {$108^\circ$};
\fill[accent!70,draw=white,line width=0.8pt] (0,0)--(-18.000:1.45) arc[start angle=-18.000,end angle=-108.000,radius=1.45cm]--cycle;
\node[font=\scriptsize\bfseries,text=white] at (-63.000:0.899) {$90^\circ$};
\fill[violet!70,draw=white,line width=0.8pt] (0,0)--(-108.000:1.45) arc[start angle=-108.000,end angle=-180.000,radius=1.45cm]--cycle;
\node[font=\scriptsize\bfseries,text=white] at (-144.000:0.899) {$72^\circ$};
\fill[gold!70,draw=white,line width=0.8pt] (0,0)--(-180.000:1.45) arc[start angle=-180.000,end angle=-240.000,radius=1.45cm]--cycle;
\node[font=\scriptsize\bfseries,text=white] at (-210.000:0.899) {$60^\circ$};
\fill[navy!70,draw=white,line width=0.8pt] (0,0)--(-240.000:1.45) arc[start angle=-240.000,end angle=-270.000,radius=1.45cm]--cycle;
\node[font=\scriptsize\bfseries,text=white] at (-255.000:0.899) {$30^\circ$};
\fill[teal!70] (1.85,0.590) rectangle ++(0.22,0.22); \node[right,font=\scriptsize] at (2.10,0.700) {Platja};
\fill[accent!70] (1.85,0.240) rectangle ++(0.22,0.22); \node[right,font=\scriptsize] at (2.10,0.350) {Muntanya};
\fill[violet!70] (1.85,-0.110) rectangle ++(0.22,0.22); \node[right,font=\scriptsize] at (2.10,0.000) {Ciutat};
\fill[gold!70] (1.85,-0.460) rectangle ++(0.22,0.22); \node[right,font=\scriptsize] at (2.10,-0.350) {Poble};
\fill[navy!70] (1.85,-0.810) rectangle ++(0.22,0.22); \node[right,font=\scriptsize] at (2.10,-0.700) {Estranger};
\end{tikzpicture}\end{center}

\ex{3} Quin gràfic triaries per a representar cada variable? Justifica-ho.\begin{tasks}(1)
  \task el nombre de germans de 30 alumnes;
  \task l'alçada de 200 persones;
  \task el partit que voten 500 persones;
  \task els minuts que tarden 50 alumnes a fer els deures.
\end{tasks}

\ex{3} Tres candidats han tret 96, 92 i 98 vots. Un cartell els representa amb un diagrama de barres en què l'eix vertical comença en 90. Per què és enganyós? Dibuixa'l bé.


% =====================================================
\seccio{5}{Paràmetres de centralització}

\ex{1} Calcula la moda, la mitjana i la mediana de: \ $7,\ \allowbreak 4,\ \allowbreak 9,\ \allowbreak 7,\ \allowbreak 3,\ \allowbreak 8,\ \allowbreak 7,\ \allowbreak 5,\ \allowbreak 6$.\recorda{\textbf{Moda}: el valor més repetit. \textbf{Mitjana}: $\bar x=\dfrac{\text{suma de les dades}}{N}$. \textbf{Mediana}: el valor central de les dades \textbf{ordenades}; si n'hi ha dos, la seua mitjana.}

\ex{1} Calcula la moda, la mitjana i la mediana de: \ $12,\ \allowbreak 15,\ \allowbreak 9,\ \allowbreak 20,\ \allowbreak 14,\ \allowbreak 16,\ \allowbreak 11,\ \allowbreak 15$.

\ex{2} La taula recull quantes vegades han anat al cine en els últims tres mesos 30 alumnes. Calcula la moda, la mitjana i la mediana.\par\smallskip\begin{center}{\renewcommand{\arraystretch}{1.3}\begin{tabular}{|c|C{0.9cm}|C{0.9cm}|C{0.9cm}|C{0.9cm}|C{0.9cm}|}\hline $x_i$ & $1$ & $2$ & $3$ & $4$ & $5$\\\hline $f_i$ & $4$ & $7$ & $9$ & $6$ & $4$\\\hline\end{tabular}}\end{center}\recorda{Amb una taula: $\bar x=\dfrac{\sum x_i\cdot f_i}{N}$. Per a la mediana, mira en la columna $F_i$ on arriba la posició central.}

\ex{2} Calcula la mitjana aproximada de les alçades dels 50 alumnes de la secció 3 (exercici 16). Quin és l'interval modal? I l'interval medià?

\ex{2} Pau ha tret 6, 7 i 5 en tres exàmens. Quina nota necessita en el quart per a tindre una mitjana de $6{,}5$?\recorda{Si coneixem la mitjana i el nombre de dades, coneixem la suma: \ $\text{suma}=\bar x\cdot N$.}

\ex{2} En una classe de 25 alumnes la nota mitjana d'un examen ha sigut $6{,}4$ i en una altra de 15 alumnes, $7{,}2$. Quina és la nota mitjana dels 40 alumnes?

\ex{3} En una empresa xicoteta, quatre persones cobren 1\,100, 1\,200, 1\,250 i 1\,300\,€ al mes, i el gerent, 4\,650\,€. Calcula la mitjana i la mediana dels sous. Quin paràmetre representa millor el sou d'un treballador?

\ex{3} En aquesta taula falta una freqüència. Troba-la sabent que la mitjana és $1{,}4$.\par\smallskip\begin{center}{\renewcommand{\arraystretch}{1.3}\begin{tabular}{|c|C{0.9cm}|C{0.9cm}|C{0.9cm}|C{0.9cm}|}\hline $x_i$ & $0$ & $1$ & $2$ & $3$\\\hline $f_i$ & $4$ & $a$ & $6$ & $2$\\\hline\end{tabular}}\end{center}


% =====================================================
\seccio{6}{Paràmetres de dispersió}

\ex{1} Calcula el rang i la desviació mitjana de: \ $3,\ \allowbreak 7,\ \allowbreak 8,\ \allowbreak 10,\ \allowbreak 12$.\recorda{Rang $=x_{\max}-x_{\min}$. \ $DM=\dfrac{\sum|x_i-\bar x|}{N}$.}

\ex{1} Calcula la variància i la desviació típica de: \ $2,\ \allowbreak 4,\ \allowbreak 5,\ \allowbreak 7,\ \allowbreak 7$.\recorda{$\sigma^2=\dfrac{\sum(x_i-\bar x)^2}{N}$ \ o bé \ $\sigma^2=\dfrac{\sum x_i^2}{N}-\bar x^{\,2}$. \ $\sigma=\sqrt{\sigma^2}$.}

\ex{2} Calcula la mitjana, la variància i la desviació típica del nombre de germans de 20 alumnes:\par\smallskip\begin{center}{\renewcommand{\arraystretch}{1.3}\begin{tabular}{|c|C{0.9cm}|C{0.9cm}|C{0.9cm}|C{0.9cm}|C{0.9cm}|}\hline $x_i$ & $0$ & $1$ & $2$ & $3$ & $4$\\\hline $f_i$ & $5$ & $8$ & $4$ & $2$ & $1$\\\hline\end{tabular}}\end{center}

\ex{2} Dos jugadors de bàsquet han fet aquests punts en cinc partits. A: $6,\ 7,\ 7,\ 8,\ 7$; \ B: $10,\ 4,\ 9,\ 5,\ 7$. Calcula la mitjana i la desviació típica de cadascun. Qui és més regular?

\ex{2} L'alçada d'un grup de persones té $\bar x=160$\,cm i $\sigma=8$\,cm. El pes d'un grup de gossos té $\bar x=40$\,kg i $\sigma=4$\,kg. Quin grup té més dispersió relativa?\recorda{Per a comparar dispersions amb mitjanes o unitats diferents: $CV=\dfrac{\sigma}{\bar x}$ (en \%).}

\ex{3} Calcula la desviació típica i el coeficient de variació de les alçades dels 50 alumnes (secció 3, exercici 16).

\ex{3} Les dades $2,\ 4,\ 5,\ 7,\ 7$ tenen $\bar x=5$ i $\sigma\approx1{,}90$. Sense tornar a fer tots els càlculs, troba la mitjana, la desviació típica i el coeficient de variació si:\begin{tasks}(2)
  \task sumem 5 a totes les dades;
  \task multipliquem totes les dades per 2.
\end{tasks}

\ex{3} Dues màquines omplin botelles d'un litre. La A ompli de mitjana 1\,002\,ml amb $\sigma=3$\,ml i la B, 998\,ml amb $\sigma=8$\,ml. Quina màquina funciona millor? Raona-ho amb el coeficient de variació.


% =====================================================
\seccio{7}{Paràmetres de posició: quartils i percentils}

\ex{1} Calcula els quartils i el rang interquartílic de: \ $4,\ \allowbreak 8,\ \allowbreak 15,\ \allowbreak 16,\ \allowbreak 23,\ \allowbreak 42,\ \allowbreak 7,\ \allowbreak 10,\ \allowbreak 12,\ \allowbreak 19,\ \allowbreak 30$.\recorda{Ordena les dades i calcula la posició $\frac{k\cdot N}{100}$ ($Q_1$: $\frac N4$; $Q_3$: $\frac{3N}{4}$). Si no és entera, agafa la dada de la posició següent; si és entera, fes la mitjana d'eixa dada i la següent.}

\ex{1} Explica què vol dir cada frase:\begin{tasks}(1)
  \task Un nadó està en el percentil 70 de pes.
  \task La nota de Lluís és el primer quartil de la classe.
\end{tasks}

\ex{2} Les edats dels 12 membres d'un club són: \ $21,\ \allowbreak 25,\ \allowbreak 18,\ \allowbreak 30,\ \allowbreak 27,\ \allowbreak 22,\ \allowbreak 19,\ \allowbreak 26,\ \allowbreak 24,\ \allowbreak 28,\ \allowbreak 20,\ \allowbreak 23$. \ Calcula els quartils i els percentils 10 i 90.

\ex{2} La taula recull els gols marcats per un equip en 30 partits. Calcula els quartils i els percentils 10 i 90.\par\smallskip\begin{center}{\renewcommand{\arraystretch}{1.3}\begin{tabular}{|c|C{0.9cm}|C{0.9cm}|C{0.9cm}|C{0.9cm}|C{0.9cm}|C{0.9cm}|}\hline $x_i$ & $0$ & $1$ & $2$ & $3$ & $4$ & $5$\\\hline $f_i$ & $3$ & $8$ & $10$ & $6$ & $2$ & $1$\\\hline\end{tabular}}\end{center}\recorda{En una taula, busca el primer valor amb $F_i\geq\frac{k\cdot N}{100}$. Si una $F_i$ coincideix exactament, fes la mitjana d'eixe valor i el següent.}

\ex{2} Amb la taula de l'exercici anterior, a quin percentil correspon, aproximadament, marcar 2 gols? Quants gols cal marcar per a estar en el quartil 3 o per damunt?

\ex{3} En quin interval estan el primer quartil, el tercer quartil i el percentil 90 de les alçades dels 50 alumnes de la secció 3 (exercici 16)?

\ex{3} En unes oposicions amb 250 aspirants, només aproven els que tenen una nota igual o superior al percentil 80. Quants aspirants aproven, aproximadament? Si Marta ha tret una nota que és el percentil 75, aprova?


% =====================================================
\seccio{8}{Diagrama de caixes}

\ex{1} Dibuixa el diagrama de caixes de les dades de l'exercici 41 (secció 7).\recorda{Cinc nombres: mínim, $Q_1$, $Me$, $Q_3$ i màxim. La \textbf{caixa} va de $Q_1$ a $Q_3$ amb una línia en la mediana; els \textbf{bigotis} arriben al mínim i al màxim.}

\ex{1} Aquest diagrama de caixes representa les edats de 80 persones que han anat a un concert. Troba el rang i el rang interquartílic. Quantes persones tenen més de 50 anys? I entre 35 i 70 anys?\par\smallskip\begin{center}\begin{tikzpicture}[baseline=(current bounding box.center),line cap=round]
\draw[muted!30,line width=0.3pt] (0.000,0)--(0.000,0.950); \draw[muted] (0.000,0.06)--(0.000,-0.06) node[below,font=\tiny,text=muted,inner sep=1.5pt]{$15$};
\draw[muted!30,line width=0.3pt] (0.667,0)--(0.667,0.950); \draw[muted] (0.667,0.06)--(0.667,-0.06) node[below,font=\tiny,text=muted,inner sep=1.5pt]{$20$};
\draw[muted!30,line width=0.3pt] (1.333,0)--(1.333,0.950); \draw[muted] (1.333,0.06)--(1.333,-0.06) node[below,font=\tiny,text=muted,inner sep=1.5pt]{$25$};
\draw[muted!30,line width=0.3pt] (2.000,0)--(2.000,0.950); \draw[muted] (2.000,0.06)--(2.000,-0.06) node[below,font=\tiny,text=muted,inner sep=1.5pt]{$30$};
\draw[muted!30,line width=0.3pt] (2.667,0)--(2.667,0.950); \draw[muted] (2.667,0.06)--(2.667,-0.06) node[below,font=\tiny,text=muted,inner sep=1.5pt]{$35$};
\draw[muted!30,line width=0.3pt] (3.333,0)--(3.333,0.950); \draw[muted] (3.333,0.06)--(3.333,-0.06) node[below,font=\tiny,text=muted,inner sep=1.5pt]{$40$};
\draw[muted!30,line width=0.3pt] (4.000,0)--(4.000,0.950); \draw[muted] (4.000,0.06)--(4.000,-0.06) node[below,font=\tiny,text=muted,inner sep=1.5pt]{$45$};
\draw[muted!30,line width=0.3pt] (4.667,0)--(4.667,0.950); \draw[muted] (4.667,0.06)--(4.667,-0.06) node[below,font=\tiny,text=muted,inner sep=1.5pt]{$50$};
\draw[muted!30,line width=0.3pt] (5.333,0)--(5.333,0.950); \draw[muted] (5.333,0.06)--(5.333,-0.06) node[below,font=\tiny,text=muted,inner sep=1.5pt]{$55$};
\draw[muted!30,line width=0.3pt] (6.000,0)--(6.000,0.950); \draw[muted] (6.000,0.06)--(6.000,-0.06) node[below,font=\tiny,text=muted,inner sep=1.5pt]{$60$};
\draw[muted!30,line width=0.3pt] (6.667,0)--(6.667,0.950); \draw[muted] (6.667,0.06)--(6.667,-0.06) node[below,font=\tiny,text=muted,inner sep=1.5pt]{$65$};
\draw[muted!30,line width=0.3pt] (7.333,0)--(7.333,0.950); \draw[muted] (7.333,0.06)--(7.333,-0.06) node[below,font=\tiny,text=muted,inner sep=1.5pt]{$70$};
\draw[muted!30,line width=0.3pt] (8.000,0)--(8.000,0.950); \draw[muted] (8.000,0.06)--(8.000,-0.06) node[below,font=\tiny,text=muted,inner sep=1.5pt]{$75$};
\draw[muted,line width=0.6pt] (0,0)--(8.000,0);
\node[right,font=\scriptsize,text=navy] at (8.100,0) {anys};
\draw[thick] (0.667,0.550)--(2.667,0.550) (4.667,0.550)--(7.333,0.550);
\draw[thick] (0.667,0.420)--(0.667,0.680) (7.333,0.420)--(7.333,0.680);
\draw[teal,thick,fill=teal!25] (2.667,0.330) rectangle (4.667,0.770);
\draw[accent,line width=1.6pt] (4.000,0.330)--(4.000,0.770);
\end{tikzpicture}\end{center}

\ex{2} Les notes d'un examen de dues classes es resumeixen en aquests diagrames de caixes. Respon i raona:\par\smallskip\begin{center}\begin{tikzpicture}[baseline=(current bounding box.center),line cap=round]
\draw[muted!30,line width=0.3pt] (0.000,0)--(0.000,1.800); \draw[muted] (0.000,0.06)--(0.000,-0.06) node[below,font=\tiny,text=muted,inner sep=1.5pt]{$0$};
\draw[muted!30,line width=0.3pt] (0.800,0)--(0.800,1.800); \draw[muted] (0.800,0.06)--(0.800,-0.06) node[below,font=\tiny,text=muted,inner sep=1.5pt]{$1$};
\draw[muted!30,line width=0.3pt] (1.600,0)--(1.600,1.800); \draw[muted] (1.600,0.06)--(1.600,-0.06) node[below,font=\tiny,text=muted,inner sep=1.5pt]{$2$};
\draw[muted!30,line width=0.3pt] (2.400,0)--(2.400,1.800); \draw[muted] (2.400,0.06)--(2.400,-0.06) node[below,font=\tiny,text=muted,inner sep=1.5pt]{$3$};
\draw[muted!30,line width=0.3pt] (3.200,0)--(3.200,1.800); \draw[muted] (3.200,0.06)--(3.200,-0.06) node[below,font=\tiny,text=muted,inner sep=1.5pt]{$4$};
\draw[muted!30,line width=0.3pt] (4.000,0)--(4.000,1.800); \draw[muted] (4.000,0.06)--(4.000,-0.06) node[below,font=\tiny,text=muted,inner sep=1.5pt]{$5$};
\draw[muted!30,line width=0.3pt] (4.800,0)--(4.800,1.800); \draw[muted] (4.800,0.06)--(4.800,-0.06) node[below,font=\tiny,text=muted,inner sep=1.5pt]{$6$};
\draw[muted!30,line width=0.3pt] (5.600,0)--(5.600,1.800); \draw[muted] (5.600,0.06)--(5.600,-0.06) node[below,font=\tiny,text=muted,inner sep=1.5pt]{$7$};
\draw[muted!30,line width=0.3pt] (6.400,0)--(6.400,1.800); \draw[muted] (6.400,0.06)--(6.400,-0.06) node[below,font=\tiny,text=muted,inner sep=1.5pt]{$8$};
\draw[muted!30,line width=0.3pt] (7.200,0)--(7.200,1.800); \draw[muted] (7.200,0.06)--(7.200,-0.06) node[below,font=\tiny,text=muted,inner sep=1.5pt]{$9$};
\draw[muted!30,line width=0.3pt] (8.000,0)--(8.000,1.800); \draw[muted] (8.000,0.06)--(8.000,-0.06) node[below,font=\tiny,text=muted,inner sep=1.5pt]{$10$};
\draw[muted,line width=0.6pt] (0,0)--(8.000,0);
\node[right,font=\scriptsize,text=navy] at (8.100,0) {nota};
\draw[thick] (2.400,1.400)--(4.000,1.400) (5.600,1.400)--(7.200,1.400);
\draw[thick] (2.400,1.270)--(2.400,1.530) (7.200,1.270)--(7.200,1.530);
\draw[teal,thick,fill=teal!25] (4.000,1.180) rectangle (5.600,1.620);
\draw[accent,line width=1.6pt] (4.800,1.180)--(4.800,1.620);
\node[left,font=\small\bfseries] at (-0.15,1.400) {A};
\draw[thick] (0.800,0.550)--(3.200,0.550) (6.400,0.550)--(8.000,0.550);
\draw[thick] (0.800,0.420)--(0.800,0.680) (8.000,0.420)--(8.000,0.680);
\draw[accent,thick,fill=accent!25] (3.200,0.330) rectangle (6.400,0.770);
\draw[accent,line width=1.6pt] (4.800,0.330)--(4.800,0.770);
\node[left,font=\small\bfseries] at (-0.15,0.550) {B};
\end{tikzpicture}\end{center}\begin{tasks}(1)
  \task Quina classe té la mediana més alta?
  \task Quina classe té les notes més homogènies?
  \task En quina classe ha aprovat (5 o més) un percentatge més gran d'alumnes?
\end{tasks}

\ex{2} Dibuixa el diagrama de caixes dels gols de l'exercici 44 (secció 7). Què indica que la mitja caixa de la dreta i la de l'esquerra siguen iguals?

\ex{2} D'unes dades sabem que el mínim és 12, $Q_1=18$, $Me=21$, $Q_3=30$ i el màxim, 34. Quin percentatge de les dades és:\begin{tasks}(4)
  \task menor que 21?
  \task entre 18 i 30?
  \task major que 18?
  \task entre 30 i 34?
\end{tasks}

\ex{3} Comprova si hi ha valors atípics en les dades $3,\ \allowbreak 21,\ \allowbreak 22,\ \allowbreak 24,\ \allowbreak 25,\ \allowbreak 26,\ \allowbreak 27,\ \allowbreak 28,\ \allowbreak 30$. Dibuixa'n el diagrama de caixes i calcula'n la mitjana i la mediana. Quin paràmetre representa millor les dades?\recorda{Una dada és \textbf{atípica} si és menor que $Q_1-1{,}5\cdot RI$ o major que $Q_3+1{,}5\cdot RI$.}

\ex{3} Inventa vuit dades que tinguen aquest resum: mínim 1, $Q_1=3$, $Me=5$, $Q_3=8$, màxim 12.

\vfill
\begin{center}\color{muted}\small\faCheckCircle[regular]\ \ Ànim! Comprova els teus resultats amb les eines interactives dels apunts de la unitat.\end{center}

\end{document}
